Supersymmetric Yang--Mills theories in signatures with several times are built by restricting the gaugino to the kernel of $k$ mutually orthogonal null vectors and letting the supersymmetry close on translations contracted with their product. We prove that every theory of this kind, defined by five axioms that Nishino's construction satisfies, has residual dimension $s+t-2k$ equal to ten (or to two, a degenerate case), and that its null vectors can be real only if $k\le\min(s,t)$. The cubic fermion terms reduce to a Fierz identity on the residual spinors, and that identity forces the dimension of the spinors to be $2(d-2)$ on a $d$-dimensional vector space, which happens only for $d=3,4,6,10$. With Majorana--Weyl spinors the signature must therefore satisfy $s-t\in\{8,0,-8\}$. There is no such theory in $(9+8n)+1$, $(10+8n)+2$, $(11+8n)+3$ or $(12+8n)+4$ for $n\ge1$, although it has been claimed that the existence proofs for the signatures $(9+m)+(1+m)$ carry over to these; the signatures $27+3$ and $19+3$ are among those excluded. For a single time the result is stronger: in $(9+8n)+1$ with $n\ge1$ no unitary theory has a massless vector multiplet, because every massless multiplet without tensor charges has at least $2^{2^{2+4n}}$ states.
